% ============================================================
%  Rapport de Stage -- Oracle Log RAG
%  Analyse automatique des logs Oracle Database
%  Tunisie Telecom / 2026
% ============================================================
\documentclass[12pt,a4paper]{report}

% ---------- Langue et encodage ----------
\usepackage[french]{babel}
\usepackage[utf8]{inputenc}
\usepackage[T1]{fontenc}

% ---------- Mise en page ----------
\usepackage[a4paper,left=2.5cm,right=2.5cm,top=3cm,bottom=2.5cm,headheight=60pt]{geometry}
\usepackage{graphicx}
\usepackage{fancyhdr}
\usepackage{titlesec}
\usepackage[hidelinks,colorlinks=true,linkcolor=blue,citecolor=blue,urlcolor=blue]{hyperref}
\usepackage{xcolor}
\usepackage{enumitem}
\usepackage{caption}
\usepackage{subcaption}
\usepackage{float}
\usepackage{setspace}
\usepackage{tikz}
\usetikzlibrary{shapes.geometric, arrows.meta, positioning, fit, backgrounds}
\usepackage{tabularx}
\usepackage{booktabs}
\usepackage{listings}
\usepackage{amsmath}
\usepackage{ragged2e}
\usepackage{lastpage}

% ---------- Couleurs ----------
\definecolor{ttblue}{RGB}{31,73,125}
\definecolor{titleblue}{RGB}{47,84,150}
\definecolor{codegreen}{RGB}{0,128,0}
\definecolor{codegray}{RGB}{128,128,128}
\definecolor{codeblue}{RGB}{0,0,200}
\definecolor{backcolour}{RGB}{245,245,245}

% ---------- Style des titres ----------
\titleformat{\chapter}[display]
  {\normalfont\bfseries\Large}
  {\Large\chaptertitlename\ \thechapter\ :}
  {6pt}
  {\Large}
\titlespacing*{\chapter}{0pt}{0pt}{30pt}

\titleformat{\section}
  {\normalfont\bfseries\large}{\thesection.}{0.6em}{}
\titleformat{\subsection}
  {\normalfont\bfseries\normalsize}{\thesubsection.}{0.6em}{}

% ---------- En-tete / pied de page ----------
\pagestyle{fancy}
\fancyhf{}
\renewcommand{\headrulewidth}{0.6pt}
\renewcommand{\footrulewidth}{0pt}
\fancyhead[L]{\textbf{\textcolor{ttblue}{Oracle Log RAG}}}
\fancyhead[R]{\textit{Tunisie Telecom --- Stage 2026}}
\fancyfoot[C]{\thepage\ / \pageref{LastPage}}

\fancypagestyle{plain}{%
  \fancyhf{}
  \fancyhead[L]{\textbf{\textcolor{ttblue}{Oracle Log RAG}}}
  \fancyhead[R]{\textit{Tunisie Telecom --- Stage 2026}}
  \fancyfoot[C]{\thepage\ / \pageref{LastPage}}
  \renewcommand{\headrulewidth}{0.6pt}
}

% ---------- Listings ----------
\lstdefinestyle{codestyle}{
  backgroundcolor=\color{backcolour},
  commentstyle=\color{codegreen},
  keywordstyle=\color{codeblue}\bfseries,
  numberstyle=\tiny\color{codegray},
  stringstyle=\color{orange},
  basicstyle=\ttfamily\small,
  breakatwhitespace=false,
  breaklines=true,
  keepspaces=true,
  numbers=left,
  numbersep=5pt,
  showspaces=false,
  frame=single,
  framerule=0.4pt,
  rulecolor=\color{ttblue!40},
  columns=fullflexible
}
\lstset{style=codestyle}

% ---------- Commandes utilitaires ----------
\newcommand{\code}[1]{\texttt{\small #1}}
\newcommand{\fichier}[1]{\texttt{\small\textcolor{ttblue}{#1}}}
\newcommand{\ora}[1]{\texttt{\textbf{#1}}}

% Placeholder figure (a remplacer par \includegraphics une fois les images disponibles)
\newcommand{\missingfig}[2][0.75\textwidth]{%
  \begin{figure}[H]
  \centering
  \fbox{\parbox{#1}{\centering\vspace{1.4cm}
    \textcolor{gray}{\large\textit{[Image : #2]}}\vspace{1.4cm}}}
  \caption{#2}
  \end{figure}
}

% ============================================================
\begin{document}
% ============================================================

% ============================================================
% PAGE DE GARDE
% ============================================================
\begin{titlepage}
\thispagestyle{empty}
\begin{center}

\vspace*{0.5cm}

{\Huge\bfseries\textcolor{ttblue}{RAPPORT DE STAGE}}\\[0.5cm]
{\Large\bfseries Analyse Automatique des Logs Oracle Database}\\[0.2cm]
{\large par Génération Augmentée par Récupération (RAG)}

\vspace{1.5cm}
\hrule height 1.5pt
\vspace{0.5cm}

\begin{tabular}{ll}
\textbf{Élaboré par :}       & Rayen \\[0.3cm]
\textbf{Société d'accueil :} & Tunisie Telecom \\[0.3cm]
\textbf{Département :}       & Administration Systèmes et Bases de Données \\[0.3cm]
\textbf{Année :}             & 2026 \\
\end{tabular}

\vspace{0.5cm}
\hrule height 1.5pt
\vspace{1.5cm}

{\large\textbf{Résumé}}\\[0.4cm]
\begin{minipage}{0.88\textwidth}
\justifying
Ce rapport présente le développement d'un système de triage automatique
des fichiers de logs Oracle Database pour Tunisie Telecom.
Le système combine une approche RAG (\emph{Retrieval-Augmented Generation})
avec un modèle de langage fine-tuné localement afin d'identifier, expliquer
et proposer des corrections aux erreurs Oracle sans recourir à un service
cloud payant.
\end{minipage}

\vfill
{\small Version finale --- \today}
\end{center}
\end{titlepage}

% ============================================================
% REMERCIEMENTS
% ============================================================
\chapter*{Remerciements}
\addcontentsline{toc}{chapter}{Remerciements}

Je tiens tout d'abord à adresser mes sincères remerciements à l'ensemble du
personnel de Tunisie Telecom, et plus particulièrement à celui du centre
technique IT de la Kasbah, pour l'accueil chaleureux qui m'a été réservé
durant toute la durée de mon stage.

Mes remerciements les plus sincères vont à mon encadreur,
\textbf{Monsieur Ben Mohamed Ahmed}, Chef de la Subdivision Administration
Systèmes et Bases de Données, qui m'a accordé sa confiance dès le premier jour.
Sa disponibilité constante, la pertinence de ses conseils et son expertise
technique m'ont énormément apporté et ont grandement contribué à la réussite
de ce stage.

Je souhaite également exprimer ma gratitude à
\textbf{Monsieur Mustapha Charfeddine Akkez}, dont les conseils avisés et le
soutien de tous les instants m'ont accompagné tout au long de cette expérience.

Plus généralement, je remercie chaleureusement toutes les personnes qui, de
près ou de loin, m'ont fait profiter de leur savoir-faire et de leur expérience
durant ce stage.

% ============================================================
% TABLE DES MATIERES
% ============================================================
\tableofcontents
\listoffigures

% ============================================================
% INTRODUCTION GENERALE
% ============================================================
\chapter*{Introduction générale}
\addcontentsline{toc}{chapter}{Introduction générale}

Étant actuellement étudiant en informatique, et pour avoir une formation
complète en tant qu'ingénieur, je dois effectuer un stage d'initiation.

L'objectif pédagogique de ce stage est de me faire découvrir la vie
professionnelle, le monde de l'entreprise et l'environnement technique.
De ce fait, ce rapport est un document représentatif de mon stage d'initiation
que j'ai effectué au sein de la division administration systèmes et bases de
données qui se trouve à la Centrale IT de la Kasbah, durant la période allant
du 22/07/2019 jusqu'au 24/08/2019.

Les bases de données ont pris une place importante en informatique, et en
particulier dans le domaine de la gestion. Actuellement elles sont au c\oe{}ur
des entreprises. Durant ce stage je profite donc d'améliorer mes connaissances
théoriques sur les SGBDs ainsi que leurs fonctionnalités et leurs architectures.

Dans ce rapport, je vais présenter dans le premier chapitre une présentation
générale de TUNISIE TELECOM dans son environnement en appuyant sur ses
principales activités. Dans le deuxième chapitre je vais parler de la partie
théorique du stage, et finalement dans le troisième chapitre je vais présenter
la partie pratique. Le rapport se termine par une conclusion qui résume
l'ensemble de nos travaux ainsi que des perspectives.

% ============================================================
% CHAPITRE 1 : PRESENTATION DE L'ENTREPRISE
% ============================================================
\chapter{Présentation de l'entreprise}

\section{Présentation générale de l'entreprise d'accueil Tunisie Telecom}

L'Office National de Télécom est un établissement public à caractère commercial.
Il a été créé le 17/04/1995 par la chambre des députés par la loi N\textdegree95-36.
Cet Office National de Télécom est né le 01/01/1996 connu par son nom commercial
Tunisie Télécom.

Elle se compose de 24 directions régionales, de 140 Espaces TT et points de vente
et de plus de 13 mille points de vente privés. Elle emploie plus de 6000 agents.

Tunisie Télécom a ainsi pour mission d'assurer les activités relatives au domaine
de télécommunication. Il est notamment chargé de~:

\begin{itemize}[leftmargin=*]
  \item L'installation, le développement, l'entretien et l'exploitation des réseaux
        de téléphone et la transmission des données.
  \item La promotion des nouveaux services de télécommunication.
  \item L'offre de tous les services publics ou privés de télécommunication
        correspondant aux divers besoins à caractère social et économique.
  \item La participation à l'effort national d'enseignement supérieur au niveau
        du secteur de télécommunication.
  \item La promotion de la coopération dans tous les domaines de télécommunications.
\end{itemize}

L'Office National de Télécommunications Tunisie Télécom est placé sous la tutelle
du ministère de la communication~; son siège est fixé à Tunis.

\section{Organigramme de Tunisie Telecom (TT)}

\missingfig{Organigramme administratif de Tunisie Télécom}

\section{La direction centrale des systèmes d'information (DCSI)}

Cette direction s'occupe de la mise en \oe{}uvre de l'infrastructure informatique
et des réseaux sécurisés entreprise (Data center) ainsi que l'administration des
systèmes et les bases de données.

Elle comprend des différents services~:

\subsection{Service réseau informatique}

Ce service a en charge l'ensemble des réseaux de télécommunications de
l'entreprise, qui couvrent les réseaux locaux et distants~: téléphonie,
internet, intranet\ldots

Il est responsable des bonnes performances du réseau, organise et définit
les procédures d'interconnexion des équipements qui constituent le réseau
informatique~: routeurs, commutateurs (switch), etc.

\subsection{Service sécurité informatique}

Ce service vise principalement les objectifs suivants~:

\begin{itemize}[leftmargin=*]
  \item \textbf{Confidentialité}~: seules les personnes autorisées peuvent avoir
        accès aux informations qui leur sont destinées (notions de droits ou
        permissions) à travers la protection des ports. Tout accès indésirable
        doit être empêché.
  \item \textbf{Authentification}~: les utilisateurs doivent prouver leur
        identité par l'usage de code d'accès.
  \item \textbf{Sécurité des données}~: Les données mises en place sur les
        serveurs de centre Data doivent être accessibles à tout moment et
        protégées des dégâts extérieurs. Il y a ainsi des protections contre les
        coupures électriques, les risques d'incendie, l'accès de personnes
        malveillantes sur les serveurs.
\end{itemize}

\missingfig{Data Center}

Il existe plusieurs méthodes pour protéger un réseau informatique~; le logiciel
le plus connu est le pare-feu (Firewall) qui sert à surveiller et contrôler les
applications et les flux des données.

\missingfig{Firewall}

\subsection{Service administration système}

Ce service est responsable de la disponibilité des informations au sein de
l'entreprise et la veille technologique dans le périmètre technique des
matériels et logiciels de type serveur, principalement l'installation,
désinstallation, la maintenance et la supervision des systèmes
d'exploitation~: Red Hat, Ubuntu, Windows-server, CentOS.

L'une des missions les plus importantes d'un \textbf{administrateur système}
consiste à veiller à la sécurité et à la sauvegarde des données sur le réseau
complet. En cas de panne ou d'incident, il doit se montrer réactif et effectuer
les réparations nécessaires dans les plus brefs délais. Son travail ne s'arrête
pas là~: grâce à une veille technologique permanente, l'administrateur systèmes
et réseaux cherche à optimiser le système en testant de nouveaux matériels.

% ============================================================
% CHAPITRE 2 : PARTIE THEORIQUE
% ============================================================
\chapter{Partie théorique}

\section{Rôle d'administrateur de base de données}

L'Administrateur de Bases de Données intervient dans la conception,
l'administration et la maintenance des bases de données, ainsi que dans
l'assistance aux informaticiens et aux utilisateurs.

L'Administrateur de Bases de Données a ainsi pour missions principales de~:

\begin{itemize}[leftmargin=*]
  \item Prévoir les volumes de données dans l'optique de choisir les outils
        servant à construire et à exploiter la base
  \item Concevoir, tester et mettre en place des progiciels SGBD (système de
        gestion de base de données)
  \item Préconiser des bonnes pratiques à usage des équipes de développement
  \item Concevoir et définir les paramètres des bases de données
  \item Réfléchir et mettre en place le dimensionnement du serveur
  \item Organiser les systèmes de gestion de bases de données tout en tenant
        compte des paramètres de cohérence, qualité et sécurité
  \item Mettre en place des normes qualité et élaborer des tableaux de bord
        pour en assurer le suivi
  \item Veiller à la disponibilité des informations et à leur facile utilisation
  \item Assister les développeurs et les techniciens d'exploitation dans
        l'exercice de leur fonction
  \item Assurer une veille technologique et un contrôle de la ou des bases
        de données
\end{itemize}

\section{Les types de systèmes de bases de données}

\subsection{Les bases de données hiérarchiques}

Les tous premiers programmes de bases de données permettaient de structurer
l'information de façon hiérarchique~: chaque enregistrement dépendait d'un seul
enregistrement. Présenté sous forme d'arbre avec ses ramifications.

\subsection{Les bases de données réseau}

Les bases de données réseau prennent le relais de façon très satisfaisante. En
permettant les relations n-n (plusieurs parents / plusieurs enfants). D'une
structure en arbre, les bases de données deviennent des graphes.

\subsection{Les bases de données relationnelles}

C'est le type de bases que l'on connaît et que l'on pratique aujourd'hui. Basé
sur l'algèbre relationnel et les travaux de E.F. Codd, il permet de modéliser
facilement et sans grosses contraintes les systèmes du monde réel et de créer
des bases de données simples à maintenir, à faire évoluer et indépendantes de
leur support. Dans ce type de bases de données, les données sont organisées
en tables.

\subsection{Les bases de données objet}

La grande idée est ici de permettre «~d'attaquer~» la base de données de façon
transparente via ses «~objets~». Les objets sont un concept de programmation qui
simplifie la création de logiciel et apporte de nombreux atouts aux projets
informatiques importants.

\medskip
\noindent\textbf{Data Base Management System (DBMS)}~: Oracle, MySQL, Microsoft
SQL Server, PostgreSQL, DB2, Microsoft Access, SQLite\ldots

\section{Historique de l'Oracle}

Software Development Laboratories a été créé en 1977. En 1979, l'entreprise
change de nom en devenant Relational Software, Inc. (RSI) et introduit son
produit Oracle V2 comme base de données relationnelle. La version 2 ne
supportait pas les transactions mais implémentait les fonctionnalités SQL
basiques de requête et jointure. Il n'y a jamais eu de version 1, pour des
raisons de marketing, la première version a été la version 2. Celle-ci
fonctionnait uniquement sur les systèmes Digital VAX/VMS.

\section{Les versions de l'Oracle}

\begin{itemize}[leftmargin=*]
  \item En 1984, la version 4 supporte la cohérence en lecture
        (\emph{read consistency}).
  \item En 1985, la version 5 supporte les requêtes distribuées, dans le cadre
        de l'intégration du modèle client-serveur avec l'arrivée des réseaux au
        milieu des années 1980.
  \item En 1988, la version 6 supporte le PL/SQL, le verrouillage de lignes
        (\emph{row-level locking}) et les sauvegardes à chaud (\emph{hot backups},
        lorsque la base de données est ouverte). Oracle met sur le marché son ERP
        Oracle Financials basé sur la base de données relationnelle Oracle Database.
  \item En 1992, la version 7 supporte les contraintes d'intégrité, les procédures
        stockées et les déclencheurs (\emph{triggers}).
  \item En 1995, acquisition d'un puissant moteur multidimensionnel, commercialisé
        sous le nom d'Oracle Express.
  \item En 1997, la version 8 introduit le développement orienté objet, et les
        applications multimédia grâce aux services \emph{Oracle interMedia},
        renommé \emph{Oracle Multimedia} depuis la version 11g.
  \item En 1999, la version 8i d'Oracle est publiée dans le but d'affiner ses
        applications avec Internet (le \emph{i} fait référence à Internet). La
        base de données comporte nativement une machine virtuelle Java.
  \item En 2001, la version 9i ajoute 400 nouvelles fonctionnalités et permet
        de lire et d'écrire des documents XML. Elle intègre le moteur OLAP~:
        le moteur Oracle Express est dorénavant référencé au sein de l'option
        Oracle OLAP.
  \item En 2003, la version 10g supporte les expressions rationnelles. Le
        \emph{g} signifie \emph{grid}~; un des atouts marketing de la 10g est
        en effet qu'elle supporte le \emph{grid computing}.
  \item En novembre 2005, la version 10g Express Edition, complètement gratuite,
        est publiée, ainsi que la version 10g Release 2.
  \item En juillet 2007, la version 11g Linux et Windows.
  \item En septembre 2009, la version 11g Release 2 est publiée.
  \item En juillet 2013, la version 12c est publiée.
\end{itemize}

\section{Architecture de l'Oracle}

Un serveur Oracle est un système qui permet de gérer les bases de données et
qui offre un moyen de gestion des informations ouvert, complet et intégré. Un
serveur Oracle est constitué d'une instance et d'une base de données.

\missingfig{Architecture de l'Oracle}

\subsection{La base de données Oracle}

Une base de données complète est constituée d'un certain nombre de fichiers~:

\begin{itemize}[leftmargin=*]
  \item \textbf{Des fichiers de données} (datafiles ou databasefiles)~: les
        fichiers de données sont pour stocker les données (essentiellement les
        tables et leurs lignes), mais aussi les autres objets Oracle connexes
        aux tables (index, vues, synonymes, database links, procédures
        stockées, etc.)
  \item \textbf{Des fichiers journaux} (logfiles ou redologfiles)~: les fichiers
        journaux servent à mémoriser toutes les modifications (validées ou non)
        faites sur la base, à des fins de reprise en cas de perte de données
        physiques. La journalisation (assurée par le processus LGWR) est un
        mécanisme permanent et obligatoire d'Oracle, pour assurer un minimum de
        sécurité des données.
  \item \textbf{Des fichiers de contrôle} (control files)~: un fichier de
        contrôle qui spécifie le nom et l'emplacement des fichiers, le nom
        de la base.
  \item Un fichier d'initialisation ou de paramétrage (init file)
  \item Des fichiers d'archivage (optionnel) pour archiver les fichiers de
        contrôle
  \item Un fichier de paramètres (optionnel) qui stocke tous les paramètres
        de la base
  \item Des fichiers de trace pour répertorier toutes les tâches et erreurs
        effectuées
\end{itemize}

\subsection{Instance Oracle}

L'instance Oracle est constituée par~:

\subsubsection{Des processus d'arrière-plan}

Ils gèrent et appliquent les relations entre les structures physiques et les
structures mémoires. Ces processus représentent le programme qui rentre
directement en interaction avec le serveur Oracle. Ils répondent à toutes les
demandes et renvoient les résultats.

Il en existe deux catégories~:

\medskip
\noindent\textbf{Les processus d'arrière-plan obligatoires~:}

\begin{itemize}[leftmargin=*]
  \item \textbf{DBWR} (Data Base Writer)~: écrit le contenu du cache de données
        dans les fichiers de données de façon périodique en mettant à jour le
        point de restauration sur les fichiers log.
  \item \textbf{LGWR} (Log Writer)~: écrit le contenu du cache de reprise dans
        les fichiers de log toutes les trois secondes, quand le cache de reprise
        est plein au tiers et quand un processus DBWR décharge des données
        modifiées.
  \item \textbf{CKPT} (Checkpoint)~: pour assurer la synchronisation et la
        cohérence des données.
  \item \textbf{SMON} (System Monitor)~: effectue la restauration lors de reprise
        après panne, nettoie les segments temporaires, fusionne certains extents
        libres contigus.
  \item \textbf{PMON} (Process Monitor)~: pour gérer les pannes des processus
        clients.
\end{itemize}

\noindent\textbf{Les processus d'arrière-plan facultatifs~:}
ARCn, LMDn, RECO, CJQ0, LMON, Snnn, Dnnn, Pnnn, LCKn, QMNn.

\medskip
\noindent\textbf{NB~: L'archivage des redo log files}

Ce processus (processus ARCH ou «~ARCHIVER~») est optionnel et permet de faire
des restaurations les plus à jour possible. En l'absence d'archivage, on ne
pourra récupérer les données que de la dernière sauvegarde. Avec l'archivage
on récupérera ces mêmes données ainsi que les modifications qui ont été faites
entre la sauvegarde et le crash (seules les transactions en cours au moment du
crash sont perdues). L'archivage permet de garder tout l'historique des fichiers
redologs, qui sont recopiés sur le répertoire d'archivage dès qu'ils sont pleins
(Switch).

\textbf{Si on n'archive pas, les redologs sont écrasés cycliquement}, puisque
rappelons-le, ils sont utilisés de manière séquentielle et circulaire~!

Il est possible d'autoriser l'archivage automatique de manière dynamique par la
commande suivante~:

\begin{lstlisting}[language=SQL]
SQL> ALTER SYSTEM ARCHIVE LOG START;
\end{lstlisting}

\subsubsection{Des structures mémoires}

Elles se composent essentiellement de deux zones mémoires~:

\begin{itemize}[leftmargin=*]
  \item \textbf{La zone mémoire allouée à la SGA} (System Global Area)~: elle
        est allouée au démarrage de l'instance et représente une composante
        fondamentale d'une instance Oracle. Son but est d'économiser les E/S.
        Elle contient~:
        \begin{itemize}[leftmargin=*]
          \item la zone mémoire partagée par tous les utilisateurs de la base
                de données
          \item le cache de tampons de la base de données (Data Base Buffer
                cache)
          \item le tampon de journalisation (redo log Buffer) : log des
                changements récents
        \end{itemize}
  \item \textbf{La zone mémoire allouée à la PGA} (Program Global Area)~: elle
        est allouée au démarrage du processus serveur. Elle est réservée à
        chaque processus utilisateur qui se connecte à la base de données Oracle
        et libérée à la fin du processus.
\end{itemize}

\section{Backup Oracle~: sauvegarde et restauration}

La commande RMAN (Recovery MANager) permet aux DBA de gérer les opérations de
sauvegarde/restauration de manière souple et optimisée.

\subsection{Les fonctions de RMAN}

RMAN offre la possibilité de~:

\begin{itemize}[leftmargin=*]
  \item réaliser des sauvegardes globales de la base, d'espaces disque logiques
        (tablespace), de fichiers de données (datafiles), de fichiers de contrôle
        (controlfiles) et de fichiers d'archive (archivelog).
  \item éviter de sauvegarder les blocs Oracle vides, ce qui permet un gain
        significatif de volume de sauvegarde.
  \item gérer les périodes de conservation des sauvegardes.
  \item dupliquer une base de données de manière simple.
  \item éditer des rapports.
\end{itemize}

\subsection{Principe de stockage des sauvegardes}

\missingfig{Principe de stockage des sauvegardes RMAN}

\noindent\textbf{BackupPiece}~: C'est une entité physique contenant 1 ou
plusieurs fichiers de données. Un fichier peut très bien être sur plusieurs
backup pieces. Leur nombre et leur taille dépendent du paramètre
\texttt{maxpiecesize}.

\noindent\textbf{BackupSet}~: C'est une entité logique contenant un ou plusieurs
backup pieces.

% ============================================================
% CHAPITRE 3 : PARTIE PRATIQUE
% ============================================================
\chapter{Partie pratique}

Dans cette partie, on commence à découvrir un peu le rôle d'un administrateur
système et aussi bien le rôle d'un administrateur de base de données, à travers
différentes fonctionnalités.

Objectifs~:
\begin{itemize}[leftmargin=*]
  \item Installation des systèmes d'exploitation Oracle Linux et Ubuntu sous
        VMware Virtual Machine.
  \item Installation d'Oracle 11g sur Ubuntu.
  \item Création d'une base de données simple.
\end{itemize}

\section{Motivation}

\subsection{Le principe de la virtualisation}

La virtualisation est l'ensemble des technologies matérielles et/ou logicielles
qui permettent de faire fonctionner plusieurs systèmes d'exploitation et/ou
plusieurs applications sur une même machine, séparément les uns des autres,
comme s'ils fonctionnaient sur des machines physiques distinctes.

Il existe beaucoup d'intérêts, par exemple~:

\begin{itemize}[leftmargin=*]
  \item Migrer facilement d'une machine virtuelle à une autre.
  \item Économiser de l'argent sur le matériel.
  \item Développer ou tester des logiciels avec possibilité de recommencer
        sans casser le système hôte.
\end{itemize}

\subsection{Les outils de virtualisation}

\subsubsection{Oracle Virtual Machine}

Oracle VM est le serveur de virtualisation de l'entreprise Oracle fondée en 1977.
OVM est un logiciel libre et gratuit qui permet l'exploitation d'une machine
serveur virtuelle.

\subsubsection{VMware}

La société VMware Inc., fondée en 1998, propose plusieurs produits liés à la
virtualisation d'architectures x86. Dans la suite de pratique, on va utiliser
VMware Server pour ses nombreux avantages~: VMware permet de virtualiser
n'importe quel SE, il permet de virtualiser plusieurs systèmes d'exploitation
sans faire de multiboot et avoir une seule configuration machine. Il permet
aussi de configurer la quantité de mémoire, l'espace disque dur et l'accès ou
non au USB, réseau, carte son, etc.

\subsection{Les systèmes d'exploitation (SE) à base Linux}

\subsubsection{Oracle Linux}

Oracle Linux est un système d'exploitation Open Source. Il convient tant à des
applications Oracle qu'à une utilisation générale. Il est compatible avec
Red Hat Enterprise Linux. À part le fait que le téléchargement d'Oracle Linux
est gratuit, ce dernier représente un système éprouvé et optimisé~: Oracle
investit énormément sur ce système. En interne, il est testé pas moins de
128\,000 heures par jour.

\subsubsection{Ubuntu}

Ubuntu est un système d'exploitation GNU/Linux basé sur la distribution Linux.
Ubuntu se définit comme un système d'exploitation utilisé par des millions de
PC en représentant une interface simple, intuitive et sécurisée. On trouve aussi
d'autres SE à base Linux comme Fedora, CentOS, RedHat, etc.

\subsection{Base de données Oracle 11g}

Oracle Database 11g est adapté aux environnements transactionnels et décisionnels
sophistiqués. Non seulement ce SGBD améliore nettement les performances de 10g
mais aussi, et surtout, il offre des avantages tels qu'une installation simple et
rapide, et des fonctions complètes d'auto-gestion.

\section{Installation de VMware}

Les étapes d'installation en gros~:

\begin{itemize}[leftmargin=*]
  \item Télécharger l'application
        \texttt{VMware-workstation-full-14.0.0-6661328}
\end{itemize}

\missingfig{Début d'installation VMware}

\section{Installation de Ubuntu sous VMware}

Les étapes d'installation en gros~:

\begin{itemize}[leftmargin=*]
  \item Télécharger d'abord le fichier image
        \texttt{ubuntu-18.10-desktop-amd64}
\end{itemize}

\missingfig{Début d'installation et personnalisation Ubuntu}

\section{Installation de Oracle Linux 6 sous VMware}

Les étapes d'installation en gros~:

\begin{itemize}[leftmargin=*]
  \item Télécharger d'abord le fichier image
        \texttt{OracleLinux-R6-U3-Server-x86\_64-dvd}
\end{itemize}

\missingfig{Début d'installation Oracle Linux 6}
\missingfig{Gestion mémoire lors de l'installation}
\missingfig{Oracle Linux 6 installé}

\section{Installation d'Oracle 11g sur Oracle Linux 6}

Cette section contient des instructions étape-par-étape, détaillées, pour
l'installation d'Oracle 11g.

\subsection{Préparation de l'OS pour installer Oracle~: Préinstallation}

\begin{itemize}[leftmargin=*]
  \item Se connecter d'abord en tant que root.
  \item Télécharger \texttt{file1.zip} et \texttt{file2.zip} liés à
        l'installation.
  \item Obtenir l'adresse IP à partir de la commande
        \lstinline|ifconfig -a|, dans ce cas l'adresse est
        \texttt{192.168.2.133}.
\end{itemize}

\missingfig{Obtention de l'adresse IP}

\begin{itemize}[leftmargin=*]
  \item Changer le nom de la machine~: \texttt{localhost} par
        \texttt{oracle}.
\end{itemize}

\missingfig{Script pour modifier le nom de la machine}

\begin{itemize}[leftmargin=*]
  \item Lancer l'installation des packages par la commande~:
        \lstinline|yum install oracle-rdbms-server-11gR2-preinstall|
\end{itemize}

\missingfig{Installation des packages}

\begin{itemize}[leftmargin=*]
  \item Suite à la dernière commande, un compte ORACLE est créé.
        Il faut juste changer le mot de passe de ce compte.
\end{itemize}

\missingfig{Configuration du mot de passe du compte Oracle}

\begin{itemize}[leftmargin=*]
  \item Créer un sous-répertoire \texttt{soft} sous \texttt{/oracle} en
        saisissant les commandes suivantes~:

\begin{lstlisting}[language=bash]
mkdir /oracle
chown oracle:oinstall /oracle
mkdir /oracle/soft
chmod a+w soft
\end{lstlisting}

        (la dernière commande permet à tous les utilisateurs le droit d'écrire.)

  \item Installer FileZilla (\texttt{FileZilla\_3.32.0\_win64-setup.exe})
        pour copier sous \texttt{/oracle/soft} les deux fichiers \texttt{.zip}
        depuis Windows vers la machine virtuelle via la connexion par adresse IP.
\end{itemize}

\missingfig{Copie des deux fichiers sous /oracle répertoire}

\begin{itemize}[leftmargin=*]
  \item Lancer l'extraction du contenu de chaque fichier \texttt{.zip} par
        la commande~:

\begin{lstlisting}[language=bash]
unzip p13390677_112040_Linux-x86-64_1of7.zip
unzip p13390677_112040_Linux-x86-64_2of7.zip
\end{lstlisting}

  \item Une fois l'extraction achevée, supprimer les fichiers zip par~:
        \lstinline|rm *.zip| (il faut garder le répertoire \texttt{database}).
  \item Installer \texttt{Xming-6-9-0-31-setup.exe} (pour le graphique)
        puis lancer une session avec Putty en activant X11.
\end{itemize}

\missingfig{Session Xming activée dans Putty}

\begin{itemize}[leftmargin=*]
  \item Se placer sous le répertoire \texttt{database}, se connecter en tant
        que \texttt{oracle} et lancer l'installation par
        \lstinline|./runInstaller|.
\end{itemize}

\missingfig{Début d'installation Oracle 11g}

\subsection{Installation d'Oracle 11g}

\begin{itemize}[leftmargin=*]
  \item Installer \texttt{Xming-6-9-0-31-setup.exe} puis lancer Xming pour
        le graphique.
\end{itemize}

\missingfig{Setup Xming}

\begin{itemize}[leftmargin=*]
  \item Exécuter la commande~: \lstinline|vi .bash_profile|
\end{itemize}

\missingfig{Édition du script .bash\_profile}

\begin{itemize}[leftmargin=*]
  \item Exécuter le script par \lstinline|. .bash_profile|.
  \item Avec une session Xming ou X11 lancée dans Putty, lancer
        l'installation du Listener par~: \lstinline|netca|
        (configuration des services Oracle Net).
\end{itemize}

\missingfig{Configuration du Listener}

\begin{itemize}[leftmargin=*]
  \item Lancer les commandes~:

\begin{lstlisting}[language=bash]
export ORACLE_HOME=/oracle/app/product/11.2.0/dbhome_1
netmgr
\end{lstlisting}

        (Oracle Net Manager)

  \item Vérifier le statut du Listener par la commande~:
        \lstinline|lsnrctl status|
\end{itemize}

\missingfig{Le Listener est configuré et actif}

L'installation se termine avec succès.

\section{Création d'une base de données simple}

Créer la base de données en utilisant l'assistant «~DBCA~» (Assistant de
Configuration de Base de Données). Avec une session Xming autorisée, on
lance l'assistant de configuration de Base de Données (DBCA). Pour le premier
écran on trouve les types d'opération suivants~:

\begin{itemize}[leftmargin=*]
  \item \textbf{Créer une base de données}~: a pour but de stocker les
        informations de production correspondant à une application.
  \item \textbf{Configurer les options de base de données}~: plusieurs
        options de configuration sont définies au niveau de la base de données.
        Trois onglets de gestion clés contiennent ces options~: Maintenance,
        Limites et Paramètres du client.
  \item \textbf{Supprimer une base de données}~: on peut aussi utiliser DBCA
        (Data Base Configuration Assistant), qui permet la suppression d'une base.
  \item \textbf{Gérer les modèles}~: l'assistant permet de définir, créer,
        modifier et supprimer des modèles de base de données personnalisée.
\end{itemize}

On coche évidemment «~Créer une base de données~».

\missingfig{Début de la configuration de la base (DBCA)}

Pour le deuxième écran, on trouve le choix du modèle de la base de données.

\missingfig{Modèle BD généraliste ou traitement transactionnel}

Suivre les étapes de configuration~:

\missingfig{Confirmation --- la Base de données est créée avec succès}

% ============================================================
% CONCLUSION
% ============================================================
\chapter*{Conclusion}
\addcontentsline{toc}{chapter}{Conclusion}

Ce rapport est le résultat des travaux réalisés au cours de la période que j'ai
passée à l'espace technique de Tunisie Télécom de la Kasbah, en effectuant les
différentes applications informatiques et les différentes tâches réalisées par
les personnels de ce centre.

Ce stage m'a été d'une grande utilité pour plusieurs raisons. En effet, il m'a
permis d'améliorer mes connaissances, de suivre de près le fonctionnement des
différentes activités des télécommunications et de m'intégrer dans le milieu
professionnel.

Finalement, j'espère que mon travail sera à la hauteur des attentes, en
souhaitant avoir atteint la cible visée par ce rapport.

\end{document}
